Šajā apakšnodaļā sniegta uz \Gls{snn} balstītas objektu atpazīšanas pieejas praktiskā lietojuma eksperimentālā pārbaude izkliedētai \Gls{tof} sensoru sistēmai.

\Gls{tof} sensoru sistēmas \qty{360}{\degree} uztverei izmantoti divu veidu sensori, lai sasniegtu iepriekš noteiktu datu ātrumu \Gls{fps}. To galvenie raksturlielumi sniegti \ref{table:sensor_charact}. tabulā. Abiem sensoriem ir atšķirīgas priekšrocības un ierobežojumi. \textit{Monstar} sensori nodrošina plašāku \Gls{fov}, taču tiem ir īsāks darbības attālums un tie darbojas spektra 850~nm daļā, ko vairāk ietekmē saules gaismas traucējumi. Savukārt \textit{Blaze 101} sensoram ir lielāks uztveres attālums, un tas darbojas pie 940~nm, kas atbilst saules starojuma spektra minimumam, samazinot traucējumus. Tomēr tam ir šaurāks \Gls{fov}. Veidojot sistēmu, šie kompromisi ņemti vērā, lai līdzsvarotu pārklājumu, attālumu un noturību pret traucējumiem.

\begin{table}[h]
\caption{ToF kameru raksturlielumi mērījumu laikā}
\label{table:sensor_charact}
\footnotesize
\centering
\setlength{\tabcolsep}{3pt}
\begin{tabularx}{\textwidth}{l>{\raggedright\arraybackslash}X>{\raggedright\arraybackslash}X>{\raggedright\arraybackslash}X>{\raggedright\arraybackslash}Xc}
\toprule
\textbf{Sensors} & \textbf{FPS} & \textbf{Attālums (m)} & \textbf{Izšķirt-\newline spēja} & \textbf{FoV (H$^{\circ}$ $\times$ V$^{\circ}$)} & \textbf{Viļņa garums (nm)} \\
\midrule
Pmd Pico Monstar \cite{picoMonstar} & 10 & 0{,}5--6{,}0 & $352 \times 287$ & $100 \times 85$ & 850  \\
Basler Blaze 101 \cite{baslerBlaze101} & 10 & 0{,}3--10{,}0 & $640 \times 480$ & $67 \times 51$ & 940 \\
\midrule
\end{tabularx}
\end{table}

Apmācība un salīdzinājumi veikti gan ar pašu iegūtām, gan publiski pieejamām datu kopām, piemēram, \textit{KITTI} LiDAR un \textit{Pandaset} uz priekšu vērsta neskenējoša LiDAR datu kopām \cite{xiao2021}. Šīs datu kopas nodrošināja etalonu ToF sensoru datu novērtēšanai salīdzinājumā ar plaši izmantotiem, uz LiDAR balstītiem uztveres modeļiem. \textit{Pandaset} neskenējošais LiDAR bija piemērots, jo tas vairāk līdzinājās ToF sensoru datiem, ļaujot tiešāk salīdzināt datu kopas un izcelt ToF sensoru priekšrocības un ierobežojumus reālos scenārijos.


